% Context from chapters/compression.tex; for mathematical reference, not compilation.
\section{Transform Coding}
\label{sec-transform-coding}

                                                                 
\subsection{Coding}

Choose an orthonormal basis, such as a wavelet basis, and write the coefficients as
\eq{
	a_m = \dotp{f}{\psi_m} \in \RR.
}
Quantization maps each coefficient to an integer using a step size $T>0$:
\eq{
	q_m = Q_T(a_m)  \in \ZZ
	\qwhereq 
	Q_T(x) = \sign(x) \left\lfloor \frac{|x|}{T} \right\rfloor.
}
This quantizer has a zero bin of width $2T$: coefficients in $(-T,T)$ are set to zero.

Like hard thresholding~\eqref{eq-hard-thresh-approx}, quantization discards coefficients smaller than $T$ in magnitude. It also rounds the retained coefficients, introducing an additional error; see Figure~\ref{fig-thresholding-vs-quantizing}.

\myfigure{
\BookDrawing{tikz/compression-threshold-quantization}
}{ 
	Hard thresholding and quantization as functions of the input coefficient. 
}{fig-thresholding-vs-quantizing}

The integer coefficients $(q_m)_m$ are encoded in a binary stream of length $R$ bits.
Sections~\ref{subsec-support-coding} and \ref{subsec-entropic-coding} describe two approaches to encoding these integers. The goal is to minimize $R$ for an acceptable distortion.

                                                                 
\subsection{Decoding}

The decoder retrieves the integers $q_m$ from the binary file and reconstructs coefficient values using
\eql{\label{eq-de-quantization}
	\tilde a_m = \sign(q_m)  \pa{|q_m|+\frac{1}{2}} T.
}
Each nonzero symbol is reconstructed at the center of its quantization bin; zero is reconstructed as zero:
\begin{center}
\image{approximation}{.4}{quantization-bins}
\end{center}

The decoded signal or image is reconstructed as
\eql{\label{eq-decompression-approx}
	\Qq_T(f) \eqdef \sum_{m \in I_T} \tilde a_m \psi_m
	=\sum_m D_T(Q_T(\dotp{f}{\psi_m}))\psi_m,\qquad I_T=\{m:|a_m|\geq T\},
}
where $D_T(q)=T\sign(q)(|q|+1/2)$ and $\sign(0)=0$. This produces a reconstruction error $\norm{f-\Qq_T(f)}$.

Let $M=\#I_T$ and let $f_M$ retain the coefficients indexed by $I_T$ without quantizing them. This is the approximation~\eqref{eq-nonlinear-approx}, with equality at the threshold retained here. Orthogonality gives $\norm{f-f_M} \leq \norm{f-\Qq_T(f)}$. The following bound separates the approximation error from the additional error caused by rounding.

\begin{prop}\label{prop-error-compression}
	With this choice of $f_M$,
	\eql{\label{eq-error-compression}
		\norm{f-\Qq_T(f)}^2 \leq \norm{f-f_M}^2 + M T^2/4
		\qwhereq
		M = \#\enscond{m}{ \tilde a_m \neq 0}.
	}
\end{prop}

\begin{proof}
The reconstruction rule~\eqref{eq-de-quantization} gives, for $|a_m|\geq T$,
\eq{
	|a_m-\tilde a_m| \leq \frac{T}{2}.
}
By orthogonality,
\eq{
	\norm{f-\Qq_T(f)}^2 = \sum_m (a_m-\tilde a_m)^2 \leq
	\sum_{ |a_m|<T } |a_m|^2 + 
	\sum_{ |a_m| \geq T } \pa{\frac{T}{2}}^2, 
}
The first sum is $\norm{f-f_M}^2$, and the second contains $M$ terms, proving the b